\documentclass[aspectratio=169,10pt]{beamer}

% Shared Beamer setup for Fall 2026 course slides.
% Clean Cal Poly-inspired theme: no custom class files or uncommon packages.
\mode<presentation>{
  \usetheme{default}
  \usefonttheme{professionalfonts}
}

\definecolor{CalPolyGreen}{HTML}{154734}
\definecolor{CalPolyGold}{HTML}{BD8B13}
\definecolor{CalPolySage}{HTML}{7FA28F}
\definecolor{CalPolyMint}{HTML}{DDF1E7}
\definecolor{CalPolySand}{HTML}{A29F86}
\definecolor{CalPolyBeige}{HTML}{EFEEDE}
\definecolor{DeckInk}{HTML}{1F2933}
\definecolor{DeckMuted}{HTML}{5F6B7A}
\definecolor{DeckSoft}{HTML}{F7F7F5}

\setbeamersize{text margin left=0.55in,text margin right=0.55in}
\setbeamercolor{normal text}{fg=DeckInk,bg=white}
\setbeamercolor{structure}{fg=CalPolyGreen}
\setbeamercolor{title}{fg=white}
\setbeamercolor{frametitle}{fg=CalPolyGreen,bg=white}
\setbeamercolor{block title}{bg=CalPolySage,fg=white}
\setbeamercolor{block body}{bg=CalPolyMint,fg=black}
\setbeamercolor{alerted text}{fg=CalPolyGold}

\setbeamerfont{title}{size=\Huge,series=\bfseries}
\setbeamerfont{frametitle}{size=\Large,series=\bfseries}
\setbeamerfont{block title}{size=\normalsize,series=\bfseries}

\setbeamertemplate{navigation symbols}{}
\setbeamertemplate{itemize item}{\color{CalPolyGold}$\blacktriangleright$}
\setbeamertemplate{itemize subitem}{\color{CalPolyGreen}$\bullet$}
\setbeamertemplate{footline}{
  \hfill{\color{DeckMuted}\scriptsize\insertframenumber/\inserttotalframenumber}
  \hspace{0.18in}\vspace{0.12in}
}
\setbeamertemplate{frametitle}{
  \vspace{0.18in}
  \hspace{0.02in}\insertframetitle\par
  \vspace{0.08in}
  \noindent\makebox[\linewidth][l]{%
    {\color{CalPolyGold}\rule{0.55in}{2pt}}%
    {\color{CalPolyGreen}\rule{\dimexpr\linewidth-0.55in\relax}{2pt}}%
  }
}

\newcommand{\CourseNumber}{Course}
\newcommand{\CourseTitle}{Course Title}
\newcommand{\LectureNumber}{0}
\newcommand{\LectureTitle}{Lecture Title}
\newcommand{\LectureDate}{Date}
\newcommand{\CourseUnit}{Unit}

\newcommand{\coursenumber}[1]{\renewcommand{\CourseNumber}{#1}}
\newcommand{\coursetitle}[1]{\renewcommand{\CourseTitle}{#1}}
\newcommand{\lecturenumber}[1]{\renewcommand{\LectureNumber}{#1}}
\newcommand{\lecturetitle}[1]{\renewcommand{\LectureTitle}{#1}}
\newcommand{\lecturedate}[1]{\renewcommand{\LectureDate}{#1}}
\newcommand{\courseunit}[1]{\renewcommand{\CourseUnit}{#1}}

\author{Kirk Duran}
\institute{California Polytechnic State University, San Luis Obispo}
\date{\LectureDate}

\newcommand{\makedecktitle}{
  \begingroup
  \setbeamercolor{background canvas}{bg=CalPolyGreen}
  \begin{frame}[plain]
    \vfill
    \begin{center}
      {\usebeamerfont{title}\color{white}\LectureTitle\par}
      \vspace{0.18in}
      {\Large\color{white}Lecture \LectureNumber\par}
    \end{center}
    \vfill
    \noindent{\color{white}\rule{\linewidth}{1.2pt}}\par
    \vspace{0.08in}
    \noindent\colorbox{CalPolyGold}{\makebox[0.24\linewidth][c]{\color{white}\Large\CourseNumber}}
    \hspace{0.12in}{\color{white}\Large Kirk Duran}
  \end{frame}
  \endgroup
}

\newenvironment{keyidea}{\begin{block}{Key idea}}{\end{block}}
\newenvironment{activity}{\begin{block}{In class}}{\end{block}}
\newenvironment{cpdefinition}{\begin{block}{Definition}}{\end{block}}
\newenvironment{exampleblockcp}{
  \begingroup
  \setbeamercolor{block title}{bg=CalPolySand,fg=white}
  \setbeamercolor{block body}{bg=CalPolyBeige,fg=black}
  \begin{block}{Example}
}{
  \end{block}
  \endgroup
}

\newcommand{\imageplaceholder}[2][0.3in]{%
  \noindent\makebox[\linewidth][c]{%
    \fcolorbox{CalPolySage}{CalPolyBeige}{%
      \parbox[c][#1][c]{0.86\linewidth}{%
        \centering
        {\color{DeckMuted}\scriptsize Image placeholder: }%
        {\color{DeckInk}\scriptsize\ttfamily\detokenize{#2}}%
      }%
    }%
  }%
}

\newcommand{\sectionframe}[1]{
  \begingroup
  \setbeamercolor{background canvas}{bg=CalPolyGreen}
  \begin{frame}[plain]
    \vfill
    {\color{white}\LARGE\bfseries #1\par}
    \vspace{0.18in}
    {\color{CalPolyGold}\rule{0.22\paperwidth}{3pt}}
    {\color{white}\rule{0.62\paperwidth}{3pt}}
    \vfill
  \end{frame}
  \endgroup
}

\newcommand{\placeholdernote}{
  \begin{block}{Placeholder}
    Replace these bullets with the final lecture material, examples, and in-class activities.
  \end{block}
}


\usepackage[T1]{fontenc}
\usepackage{lmodern}
\usepackage{amsmath,amssymb}
\pdfstringdefDisableCommands{\def\translate#1{#1}}

\graphicspath{{../assets/lecture-26/}}

% Generated illustrations and reproducible charts live with their prompts/script.
\newcommand{\lecturevisual}[3][2.25in]{%
  \begin{center}
    \includegraphics[width=0.95\linewidth,height=#1,keepaspectratio]{#2}%
  \end{center}%
}

\setbeamersize{text margin left=0.42in,text margin right=0.42in}

\coursenumber{CS 4880}
\coursetitle{Artificial Intelligence}
\lecturenumber{26}
\lecturetitle{Learning Agents: From Experience to Prediction}
\lecturedate{October 26, 2026}
\courseunit{Learning Agents and Machine Learning}


% Narrative adapted from the instructor-supplied Copy of 19_480_Winter26.pdf:
% pp.1-5: learning, approximation, supervised targets, curve fitting, least squares;
% p.6: BMI/diabetes Python demonstration; pp.7-8: classification/regression metrics;
% pp.9-10: polynomial features and limits of smooth fits.
% The supplied California-housing example motivates the polynomial pipeline;
% we retain one fixed synthetic print-job sample to make comparison controlled.
% Primary checks:
% https://aima.cs.berkeley.edu/4th-ed/pdfs/newchap02.pdf (learning-agent architecture)
% https://www.teach.cs.toronto.edu/~csc311notes/neural_networks/expressiveness.html
% https://scikit-learn.org/stable/modules/linear_model.html
% https://scikit-learn.org/stable/modules/model_evaluation.html
% https://scikit-learn.org/stable/common_pitfalls.html
% All material is in the main 50-minute lecture; no appendix or optional slides.
% Timing: agent loop / approximation / targets 10; line + Python demo 4;
% least squares + gradient descent 10; regression evaluation 5;
% classification + evaluation 8; polynomial fits / complexity 7;
% polynomial classification 3; limits / synthesis 3. Total 50 minutes.
% Concepts transfer to the lab; all worked data and predictions differ from it.
% Reproducible diagrams/results: generate-lecture-foundations.py and
% generate-narrative-charts.py in ../assets/lecture-26/.

\begin{document}
\makedecktitle
\section{From Logical Agents to Learning Agents}

\begin{frame}{Logical Agents: Correct Relative to Their Knowledge}
  \small
  A logical agent can justify conclusions through a sound inference procedure:
  \[
    \mathrm{KB}\vdash\alpha\quad\Longrightarrow\quad\mathrm{KB}\models\alpha.
  \]
  $\mathrm{KB}$ is the knowledge base; $\vdash$ means ``derivable'' and $\models$ means ``true in every model satisfying the knowledge base.''
  \begin{itemize}
    \item This is a precise form of correctness: the conclusion follows from the stated assumptions.
    \item The designer must supply useful facts, rules, action models, and goals. Inference can derive actions rather than enumerate every behavior.
    \item Incorrect or incomplete assumptions can still produce poor real-world behavior.
  \end{itemize}
  \begin{keyidea}
    The bottleneck is specifying enough of the world. Could an agent learn some of those relationships from experience?
  \end{keyidea}
\end{frame}

\begin{frame}{From Writing Relationships to Learning Them}
  \small
  \begin{columns}[T,onlytextwidth]
    \begin{column}{0.53\textwidth}
      A robot's planner needs predictions about the consequences of its actions.
      \begin{itemize}
        \item Hand-authoring a reliable rule for every operating condition is difficult.
        \item Actions produce observations and outcomes that can become examples.
        \item A learned predictor can improve the information used by an existing planner.
      \end{itemize}
      Learning and logical reasoning can coexist: learning supplies an empirical model; reasoning uses it to choose actions.
    \end{column}
    \begin{column}{0.42\textwidth}
      \lecturevisual[1.8in]{slide2.png}{A robot records the outcomes of experience.}
    \end{column}
  \end{columns}
  \begin{keyidea}
    Specify the goal and learning framework; estimate some of the relationships from data.
  \end{keyidea}
\end{frame}

\begin{frame}{The Learning-Agent Loop}
  \lecturevisual[1.9in]{learning-agent-loop.png}{Learning-agent architecture adapted from AIMA: performance element, critic, learning element, problem generator, sensors and actuators.}
  \small
  The agent acts, receives percepts and feedback, changes its performance mechanism, and acts again. Learning is part of this feedback loop.
\end{frame}

\begin{frame}{What Each Part Learns or Does}
  \begin{itemize}
    \item \textbf{Performance element:} chooses actions using its current knowledge and predictive models.
    \item \textbf{Critic:} evaluates experience against a performance standard and supplies feedback.
    \item \textbf{Learning element:} uses feedback to modify knowledge or model parameters.
    \item \textbf{Problem generator:} suggests exploratory actions that could produce useful new experience.
  \end{itemize}
  \begin{keyidea}
    Today we study supervised prediction inside this architecture: learning an input--output relationship from labeled experience. A predictor supports an agent; it is not by itself an entire agent.
  \end{keyidea}
\end{frame}

\begin{frame}{A Model Approximates an Input--Output Relationship}
  \small
  Write an unknown numerical relationship as
  \[
    Y=g(\mathbf X)+\varepsilon,\qquad
    \widehat Y=f_{\boldsymbol\theta}(\mathbf X).
  \]
  \begin{itemize}
    \item $\mathbf X$: input features; $Y$: observed target; $g$: underlying trend.
    \item $\varepsilon$: noise or effects omitted from the inputs; $\boldsymbol\theta$: model parameters.
    \item Choose a hypothesis class $\mathcal H=\{f_{\boldsymbol\theta}\}$ and use examples to find a useful member.
  \end{itemize}
  \begin{keyidea}
    Before asking how to learn a function, ask whether the chosen family can represent it. Approximation, optimization, and generalization are different questions.
  \end{keyidea}
\end{frame}

\begin{frame}{Universal Approximation: An Existence Result}
  \small
  One standard theorem: if $g:K\to\mathbb R$ is continuous on a compact (closed and bounded) set $K\subset\mathbb R^d$, then for every $\delta>0$ there exists a sufficiently wide, single-hidden-layer sigmoid network $h$ such that
  \[
    \sup_{\mathbf x\in K}|h(\mathbf x)-g(\mathbf x)|<\delta.
  \]
  \begin{itemize}
    \item $h$ is the approximating function; $\delta$ is a desired uniform error tolerance.
    \item The theorem says suitable parameters \emph{exist}; it does not provide them or promise that finite training data will identify them.
    \item It does not say a straight line, a fixed-degree polynomial, or any fitted model approximates every relationship.
  \end{itemize}
  \begin{keyidea}
    Expressive models make learning plausible. They do not remove noise, optimization difficulty, or the need for held-out evaluation.
  \end{keyidea}
\end{frame}

\section{Define the Learning Problem}
\begin{frame}{Supervised Learning: Examples Become a Predictor}
  \lecturevisual[1.45in]{training-inference.png}{Separate paths for learning from labeled examples and predicting from new observations.}
  \small
  \[
    \widehat{\boldsymbol\theta}=\arg\min_{\boldsymbol\theta}
       \underbrace{\frac1n\sum_{i=1}^n\ell(f_{\boldsymbol\theta}(\mathbf x_i),y_i)}_{\text{training loss}},
    \qquad \widehat y_{\rm new}=f_{\widehat{\boldsymbol\theta}}(\mathbf x_{\rm new}).
  \]
  A training example is a recorded pair $(\mathbf x_i,y_i)$; a model is a fitted function; a new observation supplies inputs before its target is known. Labels come from recorded outcomes or annotation.
\end{frame}

\begin{frame}{The Same Features Can Support Different Targets}
  \small
  For a printer, inputs known before starting could be
  \[
    \mathbf x=(\text{page count},\text{ink coverage},\text{queue length},\text{temperature}).
  \]
  \begin{columns}[T,onlytextwidth]
    \begin{column}{0.47\textwidth}
      \begin{block}{Regression: one numerical value}
        Target $y^{(r)}\in\mathbb R$: completion time in minutes.\\
        Prediction $\widehat y=f_{\mathbf w}(\mathbf x)$.\\
        Example loss: $(\widehat y-y^{(r)})^2$.
      \end{block}
    \end{column}
    \begin{column}{0.49\textwidth}
      \begin{block}{Binary classification: a category}
        Target $y^{(c)}\in\{0,1\}$: completes or fails.\\
        Estimate $p\approx P(Y^{(c)}=1\mid\mathbf x)$, then assign a label using a threshold.\\
        Example loss: binary cross-entropy.
      \end{block}
    \end{column}
  \end{columns}
  Targets and losses define different questions about the same situation. A completion timestamp is unavailable before starting: using it as an input would leak outcome information.
\end{frame}

\begin{frame}{Supervised Learning: Predict Print Time}
  \small
  \begin{columns}[T,onlytextwidth]
    \begin{column}{0.49\textwidth}
      A print scheduler needs a duration estimate before accepting a job.
      \begin{itemize}
        \item $x_i$: job size in ten-page batches.
        \item $y_i$: observed completion time in minutes.
        \item $\mathcal D=\{(x_i,y_i)\}_{i=1}^{n}$: labeled training examples; here $n=12$.
        \item $\widehat y=f_{\mathbf w}(x)$: prediction from a model with parameters $\mathbf w$.
      \end{itemize}
      Training learns $\mathbf w$ from observed targets. Inference predicts a target for a new job.
    \end{column}
    \begin{column}{0.47\textwidth}
      \lecturevisual[1.95in]{case-data.png}{Twelve synthetic print jobs used throughout the regression example.}
    \end{column}
  \end{columns}
\end{frame}

\begin{frame}{Curve Fitting: Start with a Line}
  \small
  Curve fitting chooses a function whose predictions match numerical observations under a specified loss. Start with
  \[
    \widehat y=w_0+w_1x
  \]
  \begin{columns}[T,onlytextwidth]
    \begin{column}{0.48\textwidth}
      \begin{itemize}
        \item $w_0$: intercept, in minutes.
        \item $w_1$: slope, in minutes per ten-page batch.
        \item A family of lines is a \textbf{hypothesis class}: each choice of weights gives one candidate predictor.
        \item Which line should we prefer? We need a numerical objective.
      \end{itemize}
      A line assumes a constant change in predicted time per added batch.
    \end{column}
    \begin{column}{0.48\textwidth}
      \lecturevisual[1.85in]{case-candidate-lines.png}{Two candidate lines on the same training jobs.}
    \end{column}
  \end{columns}
\end{frame}

\begin{frame}[fragile]{A Python Fit: BMI and a Progression Score}
  \small
  Use the diabetes benchmark's BMI feature to predict a numerical progression score. Hold out the final 20 cases for this demonstration.
{\footnotesize
\begin{verbatim}
from sklearn.datasets import load_diabetes
from sklearn.linear_model import LinearRegression

X, y = load_diabetes(return_X_y=True)
X = X[:, [2]]                  # BMI; keep a 2D feature matrix
X_train, X_test = X[:-20], X[-20:]
y_train, y_test = y[:-20], y[-20:]
model = LinearRegression().fit(X_train, y_train)
y_pred = model.predict(X_test)
\end{verbatim}
}
  \begin{keyidea}
    \texttt{fit} estimates parameters using labeled training examples; \texttt{predict} applies the fixed model to new inputs. Test targets are not passed to \texttt{fit}.
  \end{keyidea}
\end{frame}

\begin{frame}{What Did \texttt{fit} Actually Learn?}
  \small
  \begin{columns}[T,onlytextwidth]
    \begin{column}{0.48\textwidth}
      On 422 training cases, the example learns
      \[
        \widehat y\approx152.919+938.238x,
      \]
      where $x$ is the dataset-scaled BMI feature.
      \begin{itemize}
        \item The model is the fitted relationship, not a copy of a single patient record.
        \item Points show 20 held-out outcomes; departures from the line are prediction errors.
        \item Software gave us coefficients. What objective selected them, and how can an algorithm find them?
      \end{itemize}
    \end{column}
    \begin{column}{0.48\textwidth}
      \lecturevisual[1.95in]{instructor-diabetes-fit.png}{The instructor's BMI example rerun with its exact feature and split.}
    \end{column}
  \end{columns}
\end{frame}

\section{How Fitting Works}

\begin{frame}{Fit the Line by Minimizing Squared Error}
  \small
  \begin{columns}[T,onlytextwidth]
    \begin{column}{0.52\textwidth}
      Use the residual convention predicted minus observed:
      \[
        r_i=\widehat y_i-y_i=(w_0+w_1x_i)-y_i.
      \]
      Ordinary least squares chooses weights to minimize
      \[
        L(\mathbf w)=\frac1n\sum_{i=1}^{n}r_i^2,
        \qquad \widehat{\mathbf w}=\arg\min_{\mathbf w}L(\mathbf w).
      \]
      Weights are the learned parameters; loss is the scalar objective. Squaring prevents signed errors from canceling and emphasizes large misses.
    \end{column}
    \begin{column}{0.44\textwidth}
      \lecturevisual[1.9in]{case-line-errors.png}{Vertical residual segments from the least-squares line to observations.}
    \end{column}
  \end{columns}
  \medskip
  Minimizing the sum or the mean of squared errors gives the same fitted line.
  \smallskip

  The other common convention is $e_i=y_i-\widehat y_i=-r_i$.
  Since $e_i^2=r_i^2$ and $|e_i|=|r_i|$, both conventions give the same MSE, RMSE, and MAE.
\end{frame}

\begin{frame}{The Least-Squares Solution and a Prediction}
  \small
  With $\overline x$ and $\overline y$ the training-sample means,
  \[
    \widehat w_1=\frac{\sum_i(x_i-\overline x)(y_i-\overline y)}
                         {\sum_i(x_i-\overline x)^2},
    \qquad \widehat w_0=\overline y-\widehat w_1\overline x.
  \]
  \begin{columns}[T,onlytextwidth]
    \begin{column}{0.49\textwidth}
      For our 12 jobs: $\overline x=6.5$, $\overline y\approx17.405$.
      \[
        \boxed{\widehat y\approx4.130+2.042x}
      \]
      A new 60-page job has $x=6$:
      \[
        \widehat y\approx16.38\text{ minutes}.
      \]
      The slope requires distinct training inputs. The intercept at $x=0$ is outside our observed range.
    \end{column}
    \begin{column}{0.45\textwidth}
      \lecturevisual[1.55in]{case-line.png}{The fitted straight line on the twelve training jobs.}
    \end{column}
  \end{columns}
\end{frame}

\begin{frame}{How the Gradient Comes from the Loss}
  \small
  For $r_i=w_0+w_1x_i-y_i$ and $L=\frac1n\sum_i r_i^2$, the chain rule gives
  \[
    \frac{\partial r_i^2}{\partial w_0}=2r_i,
    \qquad
    \frac{\partial r_i^2}{\partial w_1}=2r_ix_i.
  \]
  Therefore
  \[
    g_0=\frac{\partial L}{\partial w_0}=\frac2n\sum_i r_i,
    \qquad
    g_1=\frac{\partial L}{\partial w_1}=\frac2n\sum_i x_i r_i.
  \]
  \begin{itemize}
    \item Each partial derivative measures sensitivity of loss to one parameter.
    \item The gradient $\nabla L=(g_0,g_1)$ points in the direction of steepest local increase.
    \item Changing the residual convention changes derivative signs in the derivation, but not the squared-loss objective or its optimum.
  \end{itemize}
\end{frame}

\begin{frame}{Gradient Descent Updates the Whole Parameter Vector}
  \small
  Starting from $\mathbf w^{(0)}$, repeatedly evaluate the gradient and move downhill:
  \[
    \boxed{\mathbf w^{(t+1)}=\mathbf w^{(t)}-\eta\nabla L(\mathbf w^{(t)})},
    \qquad \eta>0.
  \]
  Equivalently, for a line,
  \[
    w_0^{(t+1)}=w_0^{(t)}-\eta g_0^{(t)},\qquad
    w_1^{(t+1)}=w_1^{(t)}-\eta g_1^{(t)}.
  \]
  \begin{itemize}
    \item $t$ counts iterations; $\eta$ is the learning rate; $g_j^{(t)}$ is evaluated at the current weights.
    \item Compute \emph{both} derivatives before modifying either weight: this is a simultaneous update.
    \item A small enough step decreases this smooth loss away from a stationary point. An arbitrary step size need not do so.
  \end{itemize}
  Linear least squares has a convex loss: any stationary point is a global minimum; uniqueness additionally requires sufficient independent inputs.
\end{frame}

\begin{frame}{One Worked Update on Four Different Print Jobs}
  \small
  For this arithmetic example only, use $(x,y)=(1,3),(2,5),(3,4),(4,8)$. Start at $(w_0,w_1)=(0.5,1)$ and use $\eta=0.05$.
  \begin{center}
    \renewcommand{\arraystretch}{1.18}
    \begin{tabular}{lrrrr}
      $x$ & 1 & 2 & 3 & 4 \\\hline
      Observed $y$ & 3 & 5 & 4 & 8 \\
      Predicted $\widehat y$ & 1.5 & 2.5 & 3.5 & 4.5 \\
      $r=\widehat y-y$ & $-1.5$ & $-2.5$ & $-0.5$ & $-3.5$
    \end{tabular}
  \end{center}
  \[
    g_0=\frac24(-8)=-4,\qquad g_1=\frac24(-22)=-11,
    \qquad (w_0',w_1')=(0.7,1.55).
  \]
  The updated predictions are $(2.25,3.80,5.35,6.90)$. Training MSE drops from $5.25$ to $1.25875$ min$^2$.
  \medskip

  This verifies one downhill step on known data. It does not establish future prediction accuracy.
\end{frame}

\begin{frame}{Step Size Controls the Optimization Trajectory}
  \lecturevisual[1.65in]{gradient-geometry.png}{Exact four-job squared-loss contours and learning-rate trajectories.}
  \small
  On this example, $\eta=0.05$ moves toward a minimum; $\eta=0.15$ diverges. Very small steps can take many iterations. Feature scale affects a useful step size.
  \medskip

  Gradient descent is one fitting algorithm; least-squares software can instead solve the problem with numerical linear algebra.
\end{frame}

\section{How Good Are the Predictions?}

\begin{frame}{Fitting Data and Predicting Future Experience}
  \small
  Reserve independent examples before fitting or choosing models.
  \lecturevisual[0.75in]{slide26.png}{Training fits coefficients, validation chooses settings, and test reports the frozen choice.}
  \begin{itemize}
    \item \textbf{Training:} fit weights; our synthetic case uses 12 jobs.
    \item \textbf{Validation:} compare model choices on 100 different jobs.
    \item \textbf{Test:} evaluate the frozen choice on 100 untouched jobs.
  \end{itemize}
  Our real goal is small expected loss on future examples from distribution $P$:
  \[
    R(f)=\mathbb E_{(\mathbf X,Y)\sim P}[\ell(f(\mathbf X),Y)].
  \]
  Training minimizes a sample loss; representative held-out data estimate future performance. Neither is a guarantee for every individual prediction.
\end{frame}

\begin{frame}{Regression Metrics: How Far Are Predictions Off?}
  \small
  On an evaluation set of $m$ examples, let $r_i=\widehat y_i-y_i$.
  \[
    \operatorname{MSE}=\frac1m\sum_i r_i^2,\qquad
    \operatorname{RMSE}=\sqrt{\frac1m\sum_i r_i^2},\qquad
    \operatorname{MAE}=\frac1m\sum_i|r_i|.
  \]
  \begin{center}
    \renewcommand{\arraystretch}{1.22}
    \begin{tabular}{lrrl}
      \textbf{Fitted line} & \textbf{Training} & \textbf{Validation} & \textbf{Units} \\\hline
      MSE  & 2.287 & 1.374 & minutes squared \\
      RMSE & 1.512 & 1.172 & minutes \\
      MAE  & 1.241 & 0.971 & minutes
    \end{tabular}
  \end{center}
  Smaller is better on the \emph{same} evaluation set. RMSE emphasizes large errors more than MAE; neither bounds the error of an individual prediction.
  \medskip

  A two-minute miss contributes $4$ min$^2$ to squared error but $2$ min to absolute error. A four-minute miss contributes $16$ min$^2$ versus $4$ min.
\end{frame}

\begin{frame}{Relative Performance: Baselines and $R^2$}
  \small
  A simple deployable baseline always predicts the \emph{training} mean:
  \[
    \widehat y_{\rm baseline}=\overline y_{\rm train}\approx17.405\text{ minutes}.
  \]
  The coefficient of determination uses the evaluation-set mean $\overline y_{\rm eval}$:
  \[
    R^2=1-\frac{\sum_i(y_i-\widehat y_i)^2}
                    {\sum_i(y_i-\overline y_{\rm eval})^2}.
  \]
  \begin{itemize}
    \item $R^2=1$: perfect predictions on this set.
    \item $R^2=0$: same squared error as always predicting this set's mean; $R^2<0$ is possible.
    \item The evaluation mean is a reference in the metric, \emph{not} a value we use to fit our deployed baseline.
    \item $R^2$ is unitless; pair it with RMSE or MAE to judge practical error size.
  \end{itemize}
  For the BMI demonstration: held-out MSE $2548.07$, RMSE $50.48$ score units, and $R^2\approx0.473$. These interpretations assume nonconstant targets.
\end{frame}

\section{Classification and Its Evaluation}

\begin{frame}{Change the Target: Will the Printer Fail?}
  \small
  Print time was numerical. Now the observed label is $y=1$ for a failed job and $y=0$ for a completed job. Pre-job sensor readings are inputs.
  \[
    z=w_0+\sum_{j=1}^{d}w_jx_j,\qquad
    p_{\mathbf w}(\mathbf x)=\sigma(z)=\frac1{1+e^{-z}}.
  \]
  \begin{columns}[T,onlytextwidth]
    \begin{column}{0.49\textwidth}
      \begin{itemize}
        \item $z$: unrestricted linear score; $d$: number of input features.
        \item $p$: estimate of $P(Y=1\mid\mathbf x)$, between $0$ and $1$.
        \item \textbf{Logistic regression} is a classification model despite its name.
      \end{itemize}
    \end{column}
    \begin{column}{0.45\textwidth}
      \lecturevisual[1.35in]{sigmoid-score.png}{A sigmoid maps a real-valued score to a probability estimate.}
    \end{column}
  \end{columns}
\end{frame}

\begin{frame}{Fit Class Probabilities with Log Loss}
  For a binary label $y\in\{0,1\}$ and predicted failure probability $p$,
  \[
    \ell(p,y)=-\bigl[y\log p+(1-y)\log(1-p)\bigr].
  \]
  \begin{itemize}
    \item If $y=1$, loss is $-\log p$; if $y=0$, it is $-\log(1-p)$.
    \item For a failed job: $p=0.8$ gives loss $0.223$; $p=0.2$ gives $1.609$.
    \item Fit weights by minimizing average \textbf{binary cross-entropy} (log loss), using numerical optimization; software often adds a weight penalty.
  \end{itemize}
  \begin{keyidea}
    Squared error fits numerical durations; log loss fits binary class probabilities. A confidently wrong probability incurs a large log loss.
  \end{keyidea}
  \small Logs are natural. For $x_{i0}=1$, the mean-loss gradient is $\frac{\partial L}{\partial w_j}=\frac1n\sum_i(p_i-y_i)x_{ij}$: optimize the weights using the same update framework.
\end{frame}

\begin{frame}{Evaluate Decisions with a Confusion Matrix}
  \small
  Predict failure when $p\geq\tau$, where $\tau$ is a threshold.
  A score $p=0.7$ predicts failure at $\tau=0.5$, but completion at $\tau=0.8$.
  \begin{center}
    \renewcommand{\arraystretch}{1.3}
    \begin{tabular}{l|cc}
      \textbf{100 hypothetical held-out jobs} & \textbf{Predict completion} & \textbf{Predict failure} \\\hline
      \textbf{Actual completion} & 64 true negatives & 6 false positives \\
      \textbf{Actual failure} & 9 false negatives & 21 true positives
    \end{tabular}
  \end{center}
  \begin{itemize}
    \item Rows are actual classes; columns are predictions. Failure is the positive class.
    \item A false positive triggers an unnecessary intervention; a false negative misses a failure.
    \item Lowering $\tau$ can catch more failures while triggering more false alarms.
  \end{itemize}
  Choose the threshold on validation data according to the consequences of these errors.
\end{frame}

\begin{frame}{Accuracy, Precision, Recall, and $F_1$}
  \small
  From the preceding matrix: TP=21, TN=64, FP=6, FN=9. Failure is positive.
  \[
    \text{accuracy}=\frac{\mathrm{TP}+\mathrm{TN}}{\mathrm{TP}+\mathrm{TN}+\mathrm{FP}+\mathrm{FN}}=85\%,
  \]
  \[
    P=\frac{\mathrm{TP}}{\mathrm{TP}+\mathrm{FP}}\approx77.8\%,\qquad
    R=\frac{\mathrm{TP}}{\mathrm{TP}+\mathrm{FN}}=70\%,\qquad
    F_1=\frac{2PR}{P+R}\approx73.7\%.
  \]
  \begin{itemize}
    \item \textbf{Precision:} how many predicted failures really fail?
    \item \textbf{Recall:} how many actual failures are detected? $F_1$ combines precision and recall via their harmonic mean.
    \item Predicting completion for every job gives $70\%$ accuracy but zero failure recall. Accuracy can conceal failure on a rare class.
  \end{itemize}
  State the positive class and threshold; choose metrics according to operational costs. A missed failure and a one-minute duration error have different consequences.
\end{frame}

\begin{frame}{A Class Probability Is Not a Correctness Guarantee}
  \small
  Three different evaluation questions:
  \begin{itemize}
    \item \textbf{Probability quality:} held-out log loss penalizes confident wrong predictions.
    \item \textbf{Decision quality:} a confusion matrix and precision/recall describe labels at a stated threshold.
    \item \textbf{Calibration:} among many cases assigned scores near $0.8$, about $80\%$ should actually fail if those probabilities are reliable.
  \end{itemize}
  \begin{keyidea}
    A sigmoid puts scores in $(0,1)$; it does not prove calibration. Inspect held-out outcomes rather than treating a model score as certainty.
  \end{keyidea}
  Lowering a threshold changes labels and error tradeoffs, not the underlying probability estimates. Select it on validation data.
\end{frame}

\section{The Limits of a Line and Polynomial Features}

\begin{frame}{Return to Print Time: What Does the Line Miss?}
  \small
  \begin{columns}[T,onlytextwidth]
    \begin{column}{0.47\textwidth}
      The line minimized training squared error \emph{among straight lines}.
      \begin{itemize}
        \item Positive residual: the model predicts too much time.
        \item Negative residual: actual duration exceeds the prediction.
        \item Here residuals are negative at the ends and positive in the middle.
      \end{itemize}
      This curved pattern suggests a limitation of the model family. Optimization alone cannot bend a straight line.
    \end{column}
    \begin{column}{0.49\textwidth}
      \lecturevisual[2in]{case-line-residuals.png}{The fitted line leaves a curved residual pattern.}
    \end{column}
  \end{columns}
\end{frame}

\begin{frame}{Polynomial Regression Adds Curvature}
  \small
  A degree-$k$ polynomial uses powers of the input:
  \[
    \widehat y=w_0+w_1x+w_2x^2+\cdots+w_kx^k.
  \]
  \begin{columns}[T,onlytextwidth]
    \begin{column}{0.51\textwidth}
      Fit a quadratic ($k=2$) to the \emph{same 12 jobs} by minimizing the same squared loss:
      \[
        \widehat y\approx8.028+0.372x+0.1285x^2.
      \]
      Training MSE falls from $2.287$ to $0.450$ minutes squared.
      \medskip

      The marginal change is now $w_1+2w_2x$: added batches need not add a constant amount of time.
    \end{column}
    \begin{column}{0.45\textwidth}
      \lecturevisual[1.8in]{case-quadratic.png}{A fitted line and fitted quadratic on the same training jobs.}
    \end{column}
  \end{columns}
\end{frame}

\begin{frame}[fragile]{Polynomial Features Keep the Least-Squares Framework}
  \small
  For a fixed degree $k$, construct $\boldsymbol\phi(x)=(1,x,x^2,\ldots,x^k)$:
  \[
    \widehat y=\mathbf w^{\mathsf T}\boldsymbol\phi(x),\qquad
    \widehat{\mathbf w}=\arg\min_{\mathbf w}\frac1n\sum_i
       [\mathbf w^{\mathsf T}\boldsymbol\phi(x_i)-y_i]^2.
  \]
{\footnotesize
\begin{verbatim}
from sklearn.pipeline import make_pipeline
from sklearn.preprocessing import PolynomialFeatures
from sklearn.linear_model import LinearRegression

def fit_polynomial(X_train, y_train, X_valid, degree=2):
    model = make_pipeline(PolynomialFeatures(degree, include_bias=False),
                          LinearRegression())
    model.fit(X_train, y_train)
    return model, model.predict(X_valid)
\end{verbatim}
}
  The curve is nonlinear in $x$ but linear in the learned coefficients. The degree is a hyperparameter selected using validation; coefficients are fitted on training. No new test data are used to choose degree.
\end{frame}

\begin{frame}{More Flexibility Can Fit the Sample's Noise}
  \lecturevisual[1.65in]{case-flexibility.png}{Constant, linear, quadratic and degree-nine models fitted to exactly the same twelve jobs.}
  \small
  \begin{itemize}
    \item A constant predicts the training mean. A line may \textbf{underfit}: too restricted to represent curvature.
    \item Degree 9 reduces training MSE further, from $0.450$ to $0.246$.
    \item Extra flexibility may \textbf{overfit}: follow accidental details that change in another sample.
  \end{itemize}
  With increasing degree, minimum training loss cannot rise: set extra coefficients to zero.
\end{frame}

\begin{frame}{Choose Degree Using Validation Error}
  \small
  Each candidate is fitted on the same training set and scored on the same validation set.
  \begin{center}
    \renewcommand{\arraystretch}{1.35}
    \begin{tabular}{lrr}
      \textbf{Model} & \textbf{Training MSE} & \textbf{Validation MSE} \\\hline
      Constant (degree 0; training mean) & 51.997 & 36.427 \\
      Line (degree 1) & 2.287 & 1.374 \\
      Quadratic (degree 2) & 0.450 & \textbf{0.454} \\
      Degree 9 & \textbf{0.246} & 0.701
    \end{tabular}
  \end{center}
  All errors are in minutes squared. The quadratic wins on validation, although degree 9 fits the training jobs better.
  \begin{keyidea}
    Select the quadratic on validation, then freeze it. Its untouched test RMSE is $0.614$ min versus $6.420$ min for the training-mean baseline.
  \end{keyidea}
\end{frame}

\begin{frame}{Sample Size and Noise Change the Evidence}
  \lecturevisual[1.55in]{sample-noise-comparison.png}{The same generating trend under different sample sizes and noise levels, with degree chosen using validation only.}
  \small
  \begin{itemize}
    \item More data can stabilize flexible fits; more noise makes chasing details riskier. The validation winner can change; here it stays degree 2.
    \item \textbf{Bias:} missing the trend. \textbf{Variance:} sensitivity to the training sample. Noise limits attainable accuracy.
  \end{itemize}
\end{frame}

\section{Return to Classification with Polynomial Features}

\begin{frame}{Return to Classification: Two Risky Extremes}
  \small
  Suppose both unusually cold and unusually hot pre-job temperatures increase failure risk.
  \[
    \underbrace{p(t)=\sigma(w_0+w_1t)}_{\text{linear score}}
    \qquad\longrightarrow\qquad
    \underbrace{p(t)=\sigma(w_0+w_1t+w_2t^2)}_{\text{quadratic score}}.
  \]
  Here $t$ is temperature and $y\in\{0,1\}$ is the observed failure label.
  \begin{itemize}
    \item A sigmoid of a linear score is monotonic in one input: it cannot give high failure probabilities at both ends and low probabilities in the middle.
    \item Adding $t^2$ allows that pattern. At threshold $0.5$, the boundary satisfies $w_0+w_1t+w_2t^2=0$, which can have two roots.
    \item Still fit using \textbf{binary cross-entropy}, not least squares on duration.
  \end{itemize}
\end{frame}

\begin{frame}{Polynomial Features Also Change a Classifier's Shape}
  \lecturevisual[1.7in]{case-classification.png}{Fitted linear-score and quadratic-score logistic classifiers with independent validation outcomes.}
  \small
  Both models learn from the same 200 synthetic training jobs. Dots are outcomes on 200 different validation jobs; curves are estimated probabilities.
  \medskip

  Polynomial features change the score's shape; the sigmoid, binary target, log-loss objective, and threshold rule remain the same.
\end{frame}

\begin{frame}{Evaluate the More Flexible Classifier Again}
  \small
  Compare validation log loss, then freeze the choice and threshold before using the test set.
  \begin{center}
    \renewcommand{\arraystretch}{1.25}
    \begin{tabular}{lrr}
      \textbf{Model score} & \textbf{Validation log loss} & \textbf{Validation accuracy} \\\hline
      Linear & 0.692 & 56.5\% \\
      Quadratic & 0.204 & 91.5\%
    \end{tabular}
  \end{center}
  The quadratic is selected by lower validation log loss. On 200 untouched test jobs, its log loss is 0.243 and accuracy is 94.0\% at $\tau=0.5$.
  \begin{itemize}
    \item Here the test set contains 3 false positives and 9 false negatives; accuracy alone can hide an expensive error type.
    \item Polynomial regression uses squared loss for a number; logistic regression with polynomial features uses log loss for a class probability.
    \item Degree and threshold choices belong to validation, just as degree selection did in regression.
  \end{itemize}
\end{frame}

\section{Limits of These Techniques}

\begin{frame}{When a Smooth Fit Is the Wrong Representation}
  \small
  \begin{columns}[T,onlytextwidth]
    \begin{column}{0.49\textwidth}
      Some systems have operating modes, abrupt transitions, or interactions between categorical and numerical features.
      \begin{itemize}
        \item A low-degree curve can smooth over a meaningful jump.
        \item Increasing degree may add oscillation or numerical instability rather than useful structure.
        \item Continuity prevents arbitrarily small uniform error across a genuine jump.
      \end{itemize}
    \end{column}
    \begin{column}{0.47\textwidth}
      \lecturevisual[1.75in]{smooth-fit-limits.png}{A global polynomial smooths over a genuinely discontinuous relationship.}
    \end{column}
  \end{columns}
  \medskip

  The theorem assumes continuity. A good statistical fit also does not establish causal rules; better features or a different model family may be needed.
\end{frame}

\begin{frame}{Make Evaluation Match the Intended Use}
  \small
  \begin{itemize}
    \item \textbf{Available inputs:} a completion timestamp would leak the duration. Use only information available when making the prediction.
    \item \textbf{Preprocessing:} learn scaling or imputation values on training data; reuse them unchanged on validation and test inputs.
    \item \textbf{Relevant splits:} for future jobs, consider a time-based split; to predict on new printers, keep each printer's records together.
    \item \textbf{Coverage:} our duration models saw $x\in[1,12]$. Good errors there do not validate polynomial extrapolation to much larger jobs.
    \item \textbf{Uncertainty:} scores from a finite test sample vary. Keep the test set out of repeated tuning.
  \end{itemize}
  \begin{keyidea}
    The same evaluation discipline applies to numerical predictions and categorical decisions, with metrics suited to each target.
  \end{keyidea}
\end{frame}

\begin{frame}{From Deductive Guarantees to Empirical Evidence}
  \small
  \begin{enumerate}
    \item Logical reasoning derives conclusions from an explicit knowledge base; learning can estimate relationships that are difficult to specify by hand.
    \item Labeled experience defines inputs, targets, a hypothesis class, and a loss. The agent's learning element fits the model used for future predictions.
    \item Least squares and gradient descent fit numerical models. Logistic classification changes the output interpretation and training loss.
    \item Polynomial features increase representational flexibility for both tasks, while increasing the opportunity to fit noise.
    \item Training loss measures fit; representative held-out data and task-appropriate metrics supply evidence about future usefulness.
  \end{enumerate}
  \begin{keyidea}
    Approximation asks what can be represented; optimization asks what can be fitted; generalization asks what transfers to new experience.
  \end{keyidea}
\end{frame}
\end{document}
